\section{Background and scope}
\label{sec:background}

\subsection{What we mean by context poisoning}

\emph{Context poisoning} is our term for three things happening together.
Content enters an agent's context window. The agent should not have trusted
it. The agent's behavior changes as a result. The definition says nothing
about who put the content there, and that silence is deliberate: it is what
lets both communities fit inside one term.

The term is wider than prompt injection, which is one channel and one intent.
It is wider than knowledge conflict, which is one consequence. It is narrower
than ``LLM safety'', because the damage must arrive through the context
window. Two things fall outside it. Training-time poisoning and weight attacks
do not travel through the context. Jailbreaks in which the user is the
attacker have no third-party content at all.

\subsection{Two places an answer can come from}
\label{sec:twoplaces}

A model answers from one of two sources. The first is the \emph{context}: the
document, tool output or message placed in front of it. The second is what the
model learned in training and holds in its weights. The machine learning
literature calls that \emph{parametric memory}. We use the shorter
\emph{memory} for it throughout, every result table reports it as
\texttt{pm\_acc}, and one of the defenses we evaluate, ParamMute, is named for
it.

This distinction carries the whole of Section~\ref{sec:results}.

\paragraph{Context-reliance defenses.}
A large family of published methods intervenes on exactly this distinction.
A method is a \emph{context-reliance defense} when its intervention is a
\textbf{weighting between the model's context-conditioned prediction and its
memory-conditioned prediction}. The membership test is mechanical. The method
changes \emph{how much} the model leans on retrieved text. It never asks
\emph{which} retrieved text deserves trust.

These dials are built differently and are the same dial. ACD decodes
from the no-context distribution and pulls it toward the with-context one by a
weight computed from entropy. CK-PLUG interpolates the two distributions under
a single parameter. ParamMute suppresses the feed-forward layers that carry
parametric memory, forcing the model onto the retrieved text. SpARE and SHIFT
steer along a context-faithfulness direction in activation space. COIECD
detects a conflict between the two sources and pushes toward the context. SCR
makes the same trade through prompting rather than decoding.

The definition is narrow because the definition is what carries the result.
The dial runs between \emph{believe the context} and \emph{believe memory}.
No setting on it separates a real document from a planted one, because both
are context. Section~\ref{sec:distractor} measures what turning it does and
does not achieve.

\subsection{One failure, two vocabularies}

Table~\ref{tab:vocab} sets the two vocabularies side by side. The slots are
the same. Only the names and the assumed threat model differ.

\begin{table}[t]
\centering
\caption{The same four slots, named twice. The taxonomy in
Section~\ref{sec:taxonomy} uses neutral terms so that both sides fit.}
\label{tab:vocab}
\small
\begin{tabular}{@{}p{0.26\columnwidth}p{0.31\columnwidth}p{0.31\columnwidth}@{}}
\toprule
\textbf{Slot} & \textbf{Security} & \textbf{Machine learning} \\
\midrule
The content that gets in & attack, attack technique & mechanism, degradation
mechanism \\
Who placed it & an adversary, modeled explicitly & nobody; a property of the
pipeline \\
The countermeasure & defense, mitigation & fix, method \\
What success means & the adversary fails & accuracy rises \\
Evaluation style & worst case & average case \\
\bottomrule
\end{tabular}
\end{table}

The distinction is real in one respect and empty in another. It is real because
an adaptive adversary re-optimizes against a filter and a stale retrieval index
does not. A method tuned to the average case has no reason to hold at the worst
case. It is empty because the model cannot tell the difference. A poisoned
passage and a stale passage arrive through the same channel, occupy the same
positions, and corrupt the same reasoning. The provenance differs; the token
stream does not.

Our taxonomy therefore uses \emph{mechanism} as the neutral term for both, and
records provenance in a separate \emph{intent} axis. Writing "attack" at the
top level would have made the incidental half unregisterable, because an attack
presupposes an attacker.

\subsection{The agent context as a set of inflow channels}

It helps to picture the context window as a body of water fed by a fixed set of
inflow streams. Whatever enters through any stream mixes into the same body and
affects everything downstream, whether it was dumped in deliberately or arrived
as runoff. The question this SoK answers first is simple to state: how many
streams are there, and for each one, what stands guard at its mouth?

The picture needs one correction. A pond is polluted from outside, but some of
these channels feed themselves. A poisoned agent's output becomes another
agent's input, and poisoned memory reactivates in a later session. The
phenomenon is worse than the picture: contamination can propagate to other
agents and can re-enter the same agent from storage.
